\documentclass[11pt]{article}

\usepackage[margin=1in]{geometry}
\usepackage{amsmath,amssymb,amsthm,mathtools}
\usepackage{booktabs,tabularx,array}
\usepackage{microtype}
\usepackage{enumitem}
\usepackage{xcolor}
\usepackage{xurl}
\usepackage[hidelinks]{hyperref}

% PORTFOLIO-RATIFICATION-PENDING: drafted overnight 2026-07-11 per strategy
% memo 3 (draft-now, hold-submit).  Whether this earns a standalone slot or
% becomes a section of the Paper E line is a portfolio decision reserved for
% the user.  RELEASE GATES shared with the companion papers: named authors,
% artifact manifest.

\newcommand{\ii}{\mathrm{i}}
\newcommand{\id}{\mathbf{1}}
\newcommand{\C}{\mathbb{C}}
\newcommand{\R}{\mathbb{R}}
\newcommand{\tr}{\operatorname{tr}}
\definecolor{kernelcolor}{RGB}{0,92,74}
\newcommand{\Kernel}{\textcolor{kernelcolor}{\textsc{Kernel}}}

\newtheorem{theorem}{Theorem}
\newtheorem{corollary}{Corollary}
\newtheorem{remark}{Remark}

\title{Invariant Mass as Qubit Mixedness and Null-Direction
Distinguishability}
\author{Mark Schwab}
\date{July 11, 2026 --- companion draft v1 (overnight run)}

\begin{document}
\maketitle

\begin{abstract}
For a finite family of lightlike momenta written as rank-one spinor
matrices, $P=\sum_i\psi_i\psi_i^\dagger$, the invariant mass is
$\mu=\sqrt{\det P}$.  We prove, with every finite statement checked by the
Lean~4 kernel, that this Lorentzian invariant is exactly a quantum
information quantity of the normalized momentum matrix $\rho=P/\tr P$: the
squared mass-to-trace ratio is the qubit mixedness,
$\mu^2/(\tr P)^2=\det\rho=(1-\tr\rho^2)/2$; for two spinors the rest gap is
the energy-weighted trace distance of the normalized null directions,
$\mu^2=\|\psi\|^2\|\phi\|^2-|\langle\psi,\phi\rangle|^2$, the optimal
one-shot distinguishability; and for arbitrary families the mass is an
ensemble-weighted sum of pairwise distinguishability combinations.  In the
Pauli convention $P=E\,\id+\mathbf p\cdot\boldsymbol\sigma$ this reads
$m^2/E^2=2\,(1-\tr\rho^2)$: the relativistic rest content of a momentum
ensemble is precisely the purity deficit of its direction qubit.  The
oriented Pluecker coordinate $z=\psi\wedge\phi$ refines the picture: its
modulus is the mixedness datum above, while its phase lives in the
determinant $U(1)$ of the purification freedom $P=MM^\dagger$,
$M\mapsto MR$.  At fixed $P$, that phase is invisible to observables built
only from the reduced momentum matrix and to the one-particle rest spectrum;
the companion papers prove exactly which varying-phase and
two-particle settings make it visible.  ``Mass is entropy'' would overstate
the result; the exact statement concerns purity, determinant, and pairwise
distinguishability at fixed normalization, and nothing here predicts an
absolute mass scale.
\end{abstract}

\section{Setup and the purity identity}

Let $\psi_1,\dots,\psi_N\in\C^2$ and $P=\sum_i\psi_i\psi_i^\dagger$.  Under
the Pauli map each summand is a future null momentum; $\det P$ is the
Lorentzian invariant $\mu^2$, and the classical Cauchy--Binet identity
$\det P=\sum_{i<j}|\psi_i\wedge\psi_j|^2$ is kernel-checked in the
companion derivation paper \cite{CompanionDerivation}.  Here we normalize:
$T=\tr P$ and, for $T\neq0$, $\rho=P/T$, a positive unit-trace $2\times2$
matrix --- literally a qubit density matrix built from the occupied null
directions with energy weights.

\begin{theorem}[Purity--determinant bridge]\label{thm:purity}
For any $2\times2$ complex matrix $\rho$ of unit trace,
\begin{equation}
 \det\rho=\frac{1-\tr(\rho^2)}{2}.
\end{equation}
Consequently, for $T\neq0$,
\begin{equation}
 \frac{\mu^2}{T^2}
 =\det\rho
 =\frac{1-\tr\rho^2}{2}.
\end{equation}
\end{theorem}
\noindent\Kernel{}\quad(\texttt{det\_eq\_half\_one\_sub\_purity},
\texttt{normalized\_det}.)  A single occupied direction, or several
perfectly aligned directions, give a pure state ($\tr\rho^2=1$) and zero
mass; noncollinear directions give a mixed state and a positive rest gap;
in the rest frame $\rho$ is maximally mixed.

\begin{corollary}[Pauli form]\label{cor:pauli}
Writing $P=E\,\id+\mathbf p\cdot\boldsymbol\sigma$ with real coefficients,
$\tr P=2E$ and $\det P=E^2-|\mathbf p|^2$, so for $E\neq0$,
\begin{equation}
 \frac{m^2}{E^2}=2\left(1-\tr\rho^2\right),
 \qquad \rho=\frac{P}{2E}.
\end{equation}
\end{corollary}
\noindent\Kernel{}\quad(\texttt{pauli\_trace\_and\_det}; the displayed
convention is fixed by explicit Pauli matrices in the artifact.)

\section{Mass as optimal distinguishability}

\begin{theorem}[Lagrange/trace-distance form]\label{thm:distance}
For $\psi,\phi\in\C^2$,
\begin{equation}
 |\psi\wedge\phi|^2
 =\|\psi\|^2\|\phi\|^2-|\langle\psi,\phi\rangle|^2 .
\end{equation}
\end{theorem}
\noindent\Kernel{}\quad(\texttt{wedge\_normSq\_eq}.)  For unit spinors the
right side is $1-|\langle\widehat\psi,\widehat\phi\rangle|^2$, the squared
trace distance of the corresponding pure qubit states --- the optimal
one-shot discrimination probability bias.  With the null normalization
$\|\psi\|^2=2E_1$, $\|\phi\|^2=2E_2$,
\begin{equation}
 \mu=2\sqrt{E_1E_2}\;D(\widehat\psi,\widehat\phi):
\end{equation}
after energy weighting, \emph{rest mass is how well one measurement can
tell the two null directions apart}.

\begin{theorem}[Ensemble form]\label{thm:ensemble}
For any finite family,
\begin{equation}
 \det P=\sum_{i<j}
 \Bigl(\|\psi_i\|^2\|\psi_j\|^2-|\langle\psi_i,\psi_j\rangle|^2\Bigr),
\end{equation}
so with $p_i=\|\psi_i\|^2/T$,
\begin{equation}
 \frac{\mu^2}{T^2}
 =\sum_{i<j}p_ip_j
 \Bigl(1-|\langle\widehat\psi_i,\widehat\psi_j\rangle|^2\Bigr).
\end{equation}
\end{theorem}
\noindent\Kernel{}\quad(\texttt{famMomentum\_det\_eq\_pairwise}.)
Invariant mass squared is an ensemble-weighted sum of pairwise quantum
distinguishabilities.  The maximally mixed witness (two orthogonal
equal-energy directions: $\tr\rho^2=1/2$, $\mu^2/T^2=1/4$) and the pure
control (repeated direction: $\tr\rho^2=1$, $\mu=0$) are part of the
artifact (\texttt{witness\_and\_control}).

\section{Where the phase lives, and what this paper does not claim}

Factor $P=MM^\dagger$ with $M=[\psi\ \phi]$; then $z=\det M=\psi\wedge\phi$
and the normalized amplitude $A=M/\sqrt{T}$ purifies $\rho=AA^\dagger$ with
$\det A=z/T$.  Right multiplication $M\mapsto MR$, $R\in U(2)$, preserves
$P$ while $z\mapsto z\det R$: \emph{the modulus of $z$ is the observable
mixedness datum of Theorems~\ref{thm:purity}--\ref{thm:ensemble}, while its
phase lives in the determinant $U(1)$ of the purification freedom}.  At fixed
$P$, every observable built only from the reduced momentum matrix is blind to
that purification phase, and exact unitary conjugacy preserves the one-particle
rest spectrum.  The companion papers establish exactly which additional
settings retain it: a spatially varying phase leaves derived
link data with an integer winding, and a two-particle interference
experiment separates equal-modulus fields with survival probability
$4/5$ versus $1$ \cite{CompanionDerivation}.

Four boundaries.  First, ``mass is entropy'' overstates the theorem: the
exact statements concern the determinant, purity, and pairwise
distinguishability at fixed trace normalization, not a von Neumann entropy
production principle.  Second, nothing here selects an absolute scale: the
identities are homogeneous, and the companion papers prove the natural
positive homogeneous action selects only the collapsed zero scale.  Third,
Pluecker coordinates appear widely in quantum-information treatments of
entanglement; the contribution claimed here is the exact, kernel-checked
bridge from null-momentum ensembles to Lorentzian mass, the canonical rest
operator, and operational discrimination --- not the use of Pluecker
coordinates in quantum information as such.  Fourth, the relativistic
quantum information literature has long studied how the mixedness of
\emph{reduced} states varies with the observer: reduced spin density
matrices of massive particles are not Lorentz covariant and their entropy
has no invariant meaning \cite{PeresTerno}, and boosts interchange spin and
momentum entanglement \cite{GingrichAdami}.  The present identity is
complementary rather than overlapping: it does not track mixedness across
frames but pins it, frame by frame, to the mass-to-trace ratio --- the
numerator $\mu^2$ is the Lorentz invariant, and all frame dependence is
carried by the trace normalization $T=2E$.

\appendix
\section{Machine verification}

All displayed identities are kernel-checked in Lean~4 over Mathlib in the
module \texttt{PhysicsSM/Draft/NullEdge/MassMixedness.lean}, with
build-enforced assumption footprints equal to
$\{\mathtt{propext},\mathtt{Classical.choice},\mathtt{Quot.sound}\}$, and
integrate with the project's existing trace-identity and visibility-duality
modules.  The archival release is a shared gate with the companion papers.

\begin{thebibliography}{9}

\bibitem{CompanionDerivation}
Null-Edge Formalization Project,
\emph{Null-spinor area as the rest gap of an exactly unitary Dirac walk},
companion paper, in preparation (2026).

\bibitem{PeresTerno}
A.~Peres and D.~R.~Terno,
\emph{Quantum information and relativity theory},
Rev.\ Mod.\ Phys.\ \textbf{76}, 93 (2004), arXiv:quant-ph/0212023.

\bibitem{GingrichAdami}
R.~M.~Gingrich and C.~Adami,
\emph{Quantum entanglement of moving bodies},
Phys.\ Rev.\ Lett.\ \textbf{89}, 270402 (2002), arXiv:quant-ph/0205179.

\end{thebibliography}

\end{document}
